% © 2026 Ing. David Jaroš — CC BY-NC-ND 4.0
%
% This work is licensed under a Creative Commons Attribution-NonCommercial-NoDerivatives
% 4.0 International License (CC BY-NC-ND 4.0).
%
% License History: Earlier drafts (up to v0.3) were released under CC BY 4.0.
% From v0.4 onward, all material is released under CC BY-NC-ND 4.0 to protect
% the integrity of the theoretical work during ongoing academic development.
%
% Φ_UBT from First Principles — Approaches L1–L4
% Track: CORE — α derivation — v69 baseline
% Date: 2026-03-09

\ifdefined\INCLUDEMODE
  % Being included in another document — skip preamble
\else
  \documentclass[11pt,a4paper]{article}
  \usepackage{amsmath,amssymb,amsthm}
  \usepackage{hyperref}
  \usepackage{booktabs}

  \newtheorem{theorem}{Theorem}
  \newtheorem{lemma}{Lemma}
  \newtheorem{proposition}{Proposition}
  \newtheorem{conjecture}{Conjecture}
  \theoremstyle{definition}
  \newtheorem{gap}{Gap}
  \theoremstyle{remark}
  \newtheorem{remark}{Remark}

  \title{$\Phi_{\mathrm{UBT}}$ from First Principles:\\
    \large Approaches L1 (Photon Zero-Mode),
    L2 (Atiyah--Singer),\\
    L3 (Zeta-Det), and L4 (Renormalization Finiteness)}
  \author{Ing.\ David Jaroš}
  \date{2026-03-09 (v69)}

  \begin{document}
  \maketitle
\fi

% ─────────────────────────────────────────────────────────────────────────────
\section{The Bottleneck: What is $\Phi_{\mathrm{UBT}}$ and Why is it Missing?}
\label{sec:phi_ubt_bottleneck}
% ─────────────────────────────────────────────────────────────────────────────

\textbf{Non-circularity constraint throughout this document:}
$\alpha = 1/137$ and $n^* = 137$ are \emph{never} used as inputs.
If $R_\psi$ enters a formula, it must be fixed algebraically; if it cannot,
the approach is classified as \textbf{[DEAD END]}.
All inputs are: the UBT action $\mathcal{S}[\Theta]$,
the $\mathbb{C}\otimes\mathbb{H}$ algebra ($N_{\mathrm{eff}}=12$ proved),
$\theta_H^* = 2\pi$ (Lemma~1, proved), and standard 5D heat-kernel normalization.

\subsection{The self-consistency equation (v1.1 baseline)}

The v1.1 equation for the fine-structure constant reads
\begin{equation}
  \label{eq:alpha_sc}
  \alpha \;=\;
  \frac{\theta_H^{*2}}{4\pi\,(\mathcal{I}_\psi^*)^2}
  \Bigg[
    \Phi_{\mathrm{UBT}}\sum_B C_5^{(B)}\,\mathrm{Tr}_{\mathcal{F}_B}(Q^2)
    \;+\;
    \sum_B \frac{b_B}{4\pi}
      \int_0^{\Lambda}\!\rho_B(m)\,
      \ln\!\frac{\Lambda}{m}\,dm
  \Bigg],
\end{equation}
where $\Lambda = \Phi_{\mathrm{UBT}}/Z^*$ is the UV scale.
The following quantities are \textbf{already established}:

\begin{center}
\begin{tabular}{llll}
  \toprule
  Symbol & Expression & Numerical value & Status \\
  \midrule
  $\theta_H^*$ & $2\pi n,\; n\in\mathbb{Z}\setminus\{0\}$ & $2\pi$ (canonical) & \textbf{[PROVED]} \\
  $\mathcal{I}_\psi^*$ & $\theta_H^*$ (canonical vacuum) & $2\pi$ & \textbf{[PROVED]} \\
  $\theta_H^{*2}/(4\pi(\mathcal{I}_\psi^*)^2)$ & $1/(4\pi)$ & $0.07958$ & \textbf{[PROVED]} \\
  $\mathrm{Tr}(Q^2) = N_{\mathrm{eff}}$ & $\dim_\mathbb{R}(\mathbb{C}\otimes\mathbb{H})/2$ & $12$ & \textbf{[PROVED]} \\
  $C_5$ & $1/(16\pi^2)$ & $0.006333$ & standard 5D \\
  $C_5 \times N_{\mathrm{eff}}$ & $12/(16\pi^2)$ & $0.07599$ & \textbf{[PROVED]} \\
  $\rho_B(m)$ & KK density on $S^1$ & $2R_\psi$ (leading) & derivable \\
  \midrule
  $\Phi_{\mathrm{UBT}} = \Lambda Z^*$ & pure number & $\approx 61$--$71$\,? & \textbf{[MISSING — this document]} \\
  \bottomrule
\end{tabular}
\end{center}

\subsection{Numerical constraint on $\Phi_{\mathrm{UBT}}$}

With $\theta_H^{*2}/(4\pi(\mathcal{I}_\psi^*)^2) = 1/(4\pi)$
and KK spectral density $\rho(m) = 2R_\psi$ (leading order on $S^1$),
the spectral integral evaluates to
\begin{equation}
  \frac{N_{\mathrm{eff}}}{4\pi}\int_0^\Lambda 2R_\psi\,\ln\frac{\Lambda}{m}\,dm
  \;=\;
  \frac{N_{\mathrm{eff}}}{4\pi}\cdot 2R_\psi\Lambda
  \;=\;
  \frac{N_{\mathrm{eff}}}{2\pi}\;(\text{since } R_\psi\Lambda = 1).
  \label{eq:spectral_leading}
\end{equation}
Equation~\eqref{eq:alpha_sc} then requires
\begin{equation}
  4\pi\alpha
  \;=\;
  \Phi_{\mathrm{UBT}}\cdot\frac{3}{4\pi^2}
  + \frac{N_{\mathrm{eff}}}{2\pi}
  \qquad\Longrightarrow\qquad
  \Phi_{\mathrm{UBT}}
  = \frac{4\pi^2}{3}\!\left(4\pi\alpha - \frac{N_{\mathrm{eff}}}{2\pi}\right).
  \label{eq:phi_constraint}
\end{equation}
For the canonical choice $Z^* = 1/(2\pi R_\psi)$ (Hosotani normalization)
the self-consistency equation yields $\alpha \approx 1/445$
(\textbf{too small}).
For the empirically favourable $R_\psi \approx 0.917$ the equation yields
$\alpha \approx 1/137$, but this $R_\psi$ is not yet fixed algebraically —
that is the \emph{bottleneck} addressed here.

% ─────────────────────────────────────────────────────────────────────────────
\section{Approach L1: Photon Zero-Mode Norm on $M^4\times S^1$}
\label{sec:L1_zeromode}
% ─────────────────────────────────────────────────────────────────────────────

\textbf{Status: [OPEN PROBLEM — $R_\psi$ not algebraically fixed at this level.]}

\subsection{Setup}

In the Hosotani vacuum on $M^4\times S^1_\psi$ (radius $R_\psi$) the
gauge field develops a non-trivial background along the extra dimension:
\begin{equation}
  A_\psi^{(0)} = \frac{\theta_H^*}{2\pi R_\psi} = \mathrm{const}.
  \label{eq:hosotani_bg}
\end{equation}
This is the \emph{photon zero-mode} — the unique $\psi$-independent mode
of $A_\psi$ consistent with the quantised holonomy $\theta_H^* = 2\pi$.

\subsection{Zero-mode norm $Z^*$}

The zero-mode norm is the $L^2$ norm of $A_\psi^{(0)}$ over the compact
dimension, normalised by the 4D reference volume $V_4$:
\begin{equation}
  Z^*
  = \frac{1}{V_4}\int_0^{2\pi R_\psi}\!\!\bigl|A_\psi^{(0)}\bigr|^2 d\psi
  = \frac{(2\pi R_\psi)}{V_4}\left(\frac{\theta_H^*}{2\pi R_\psi}\right)^{\!2}
  = \frac{(\theta_H^*)^2}{2\pi R_\psi\, V_4}.
  \label{eq:Z_star}
\end{equation}
For the canonical UBT choice $\theta_H^* = 2\pi$ and
$V_4 = (2\pi)^4$ (dimensionless unit volume in $\Theta$-mode normalisation):
\begin{equation}
  Z^* = \frac{(2\pi)^2}{2\pi R_\psi (2\pi)^4}
      = \frac{1}{(2\pi)^3 R_\psi}.
  \label{eq:Z_star_canonical}
\end{equation}

\subsection{$\Phi_{\mathrm{UBT}}$ from $Z^*$}

With UV cutoff $\Lambda = 1/R_\psi$ (inverse circle radius) one obtains
\begin{equation}
  \Phi_{\mathrm{UBT}}
  = \Lambda Z^*
  = \frac{1}{R_\psi}\cdot\frac{1}{(2\pi)^3 R_\psi}
  = \frac{1}{(2\pi)^3 R_\psi^2}.
  \label{eq:phi_L1}
\end{equation}
\textbf{Problem:} $\Phi_{\mathrm{UBT}}$ depends on $R_\psi$, which is not
fixed by~\eqref{eq:phi_L1} itself.

\subsection{Numerical trap: $R_\psi = \pi/\sqrt{N_{\mathrm{eff}}}$}

Substituting~\eqref{eq:phi_L1} into the self-consistency equation and
requiring $\alpha = 1/137$ determines $R_\psi$ numerically.
One finds
\begin{equation}
  R_\psi \approx 0.917,
\end{equation}
which is close (error $\approx 1\%$) to the algebraically suggestive
expression
\begin{equation}
  R_\psi^{\mathrm{alg}}
  = \frac{12\pi}{N_{\mathrm{eff}}^{3/2}}
  = \frac{N_{\mathrm{eff}}\,\pi}{N_{\mathrm{eff}}\,\sqrt{N_{\mathrm{eff}}}}
  = \frac{\pi}{\sqrt{N_{\mathrm{eff}}}}
  \overset{N_{\mathrm{eff}}=12}{=} \frac{\pi}{\sqrt{12}}
  \approx 0.907.
  \label{eq:R_psi_algebraic}
\end{equation}
The formula $R_\psi = \pi/\sqrt{N_{\mathrm{eff}}}$ has a plausible algebraic
interpretation: the $S^1_\psi$ circumference
$2\pi R_\psi = 2\pi^2/\sqrt{N_{\mathrm{eff}}}$
is the ratio of the topological period $2\pi^2$ to the root of the
$\mathbb{C}\otimes\mathbb{H}$ dimension.
However, no first-principles derivation of~\eqref{eq:R_psi_algebraic}
from $\mathcal{S}[\Theta]$ has been found.

\begin{gap}
  The algebraic origin of $R_\psi = \pi/\sqrt{N_{\mathrm{eff}}}$
  (or equivalently of the numerical $R_\psi \approx 0.917$) is unknown.
  If this formula were derivable from the UBT action, it would
  fix $\Phi_{\mathrm{UBT}}$ from equation~\eqref{eq:phi_L1} as the
  pure number
  \begin{equation}
    \Phi_{\mathrm{UBT}}
    \;=\; \frac{N_{\mathrm{eff}}}{(2\pi)^3\pi^2}
    \;\approx\; \frac{12}{(2\pi)^3\pi^2}
    \;\approx\; 0.0153,
    \label{eq:phi_L1_algebraic}
  \end{equation}
  completing the derivation.
  The $1\%$ discrepancy between $R_\psi \approx 0.917$ and
  $R_\psi^{\mathrm{alg}} \approx 0.907$ remains to be explained.
\end{gap}

\textbf{Conclusion for L1:}
The photon zero-mode approach gives $\Phi_{\mathrm{UBT}} = 1/((2\pi)^3 R_\psi^2)$
and reduces the problem to fixing $R_\psi$.
The candidate $R_\psi = \pi/\sqrt{N_{\mathrm{eff}}}$ is a
\textbf{[NUMERICAL OBSERVATION]} at $\approx 1\%$ accuracy,
not yet a proof.
Status: \textbf{[OPEN PROBLEM]}.

% ─────────────────────────────────────────────────────────────────────────────
\section{Approach L2: Atiyah--Singer Index on $\mathbb{C}\otimes\mathbb{H}$ Bundle}
\label{sec:L2_AS}
% ─────────────────────────────────────────────────────────────────────────────

\textbf{Status: [INCONCLUSIVE — index vanishes for flat base, twisted case needs
further development.]}

\subsection{Setup}

The biquaternionic field $\Theta \in \mathbb{C}\otimes\mathbb{H}$ defines
a complex vector bundle $\mathcal{E}_\Theta$ of rank $8$ (real rank 16)
over $M^4 \times S^1_\psi$.
The Hosotani background~\eqref{eq:hosotani_bg} twists this bundle through
the holonomy
\begin{equation}
  \mathrm{Hol}(\mathcal{E}_\Theta)
  = \exp\!\Bigl(\,i\oint_{S^1}A_\psi\,d\psi\Bigr)
  = e^{i\theta_H^*}
  = e^{2\pi i} = 1
  \quad(\text{for } \theta_H^* = 2\pi).
\end{equation}
Thus the Hosotani holonomy is \emph{trivial} for $\theta_H^* = 2\pi n$ with
$n\in\mathbb{Z}$, which means $\mathcal{E}_\Theta$ is topologically
equivalent to the trivial bundle in the $S^1_\psi$ direction.

\subsection{Atiyah--Singer index}

For the Dirac operator $D_\Theta$ associated with the $\Theta$-bundle
over $M^4\times S^1_\psi$:
\begin{equation}
  \mathrm{ind}(D_\Theta)
  = \int_{M^4\times S^1} \mathrm{ch}(\mathcal{E}_\Theta)\wedge\hat{A}(TM).
\end{equation}
On flat $\mathbb{R}^4\times S^1$:
\begin{itemize}
  \item $\hat{A}(T(\mathbb{R}^4\times S^1)) = 1$ (flat manifold).
  \item For the trivial bundle ($\theta_H^* = 2\pi$):
    $\mathrm{ch}(\mathcal{E}_\Theta) = \mathrm{rank}(\mathcal{E}_\Theta)
    + 0 + \cdots = 8$.
  \item $\chi(\mathbb{R}^4\times S^1) = 0$ (non-compact $\mathbb{R}^4$).
\end{itemize}
Therefore $\mathrm{ind}(D_\Theta) = 0$ for the canonical flat vacuum.

\subsection{Twisted case}

For a bundle with non-trivial first Chern class $c_1(\mathcal{E})$
(arising from a fractional holonomy $a = \theta_H^*/(2\pi) \notin\mathbb{Z}$),
one obtains a non-zero index:
\begin{equation}
  \mathrm{ind}(D_\Theta^a) = \int_{M^4\times S^1} c_1^{\mathrm{eff}}(\mathcal{E})\wedge\cdots
  \qquad (a\notin\mathbb{Z}).
\end{equation}
In this case $\Phi_{\mathrm{UBT}}$ could be identified with
$\mathrm{ind}(D_\Theta^a) / (\text{normalisation})$ as a topological
invariant.
However, since Lemma~1 establishes $\theta_H^* = 2\pi n$ (integer $n$),
the physical vacuum has trivial holonomy and $\mathrm{ind} = 0$.

\begin{remark}
  A non-trivial index could arise if the compact dimension carries torsion
  in its Pontryagin class (e.g.\ via a warped geometry or orbifold
  $S^1/\mathbb{Z}_2$).  This possibility is not excluded by Lemma~1
  and deserves further investigation.
\end{remark}

\textbf{Conclusion for L2:}
The Atiyah--Singer index vanishes for the canonical flat vacuum with
integer holonomy.
A non-trivial index relevant to $\Phi_{\mathrm{UBT}}$ would require
either a non-integer effective holonomy or a non-flat base manifold.
Status: \textbf{[INCONCLUSIVE]}.

% ─────────────────────────────────────────────────────────────────────────────
\section{Approach L3: Zeta-Regularised Determinant of $\partial_\psi + A_\psi$}
\label{sec:L3_zetadet}
% ─────────────────────────────────────────────────────────────────────────────

\textbf{Status: [KEY RESULT — gives $R_\psi$-independent pure number
$Z^* = 2\lvert\sin(\pi a)\rvert$ for $a\notin\mathbb{Z}$;
canonical vacuum $a=1$ is degenerate (zero mode) and requires separate treatment.]}

\subsection{Setup}

Consider the first-order operator on $S^1_\psi$ of radius $R_\psi$:
\begin{equation}
  \mathcal{D}_a = \partial_\psi + A_\psi,
  \qquad A_\psi = \frac{\theta_H^*}{2\pi R_\psi},
  \qquad a := \frac{\theta_H^*}{2\pi}.
  \label{eq:D_a}
\end{equation}
The eigenvalues are
$\lambda_n = (n + a)/R_\psi$, $n\in\mathbb{Z}$,
and the $\zeta$-regularised functional determinant is
$\det(\mathcal{D}_a) = \exp(-\zeta'(0;\mathcal{D}_a))$
where
$\zeta(s;\mathcal{D}_a) = \sum_{n}\lvert\lambda_n\rvert^{-s}$.

\subsection{General result for non-integer $a$}
\label{subsec:L3_noninteger}

For $a\notin\mathbb{Z}$ all eigenvalues are non-zero.
Splitting the sum into $n\geq 0$ (positive eigenvalues, using $n+a > 0$ for
large $n$) and $n\leq -1$ (negative, $\lvert n+a\rvert = -(n+a)$ for
$n < -a$) and invoking the Hurwitz zeta function
$\zeta_H(s,q)=\sum_{k=0}^\infty(k+q)^{-s}$, one obtains
\begin{equation}
  \zeta(s;\mathcal{D}_a)
  = R_\psi^{-s}\bigl[\zeta_H(s,a) + \zeta_H(s,1-a)\bigr].
  \label{eq:zeta_split}
\end{equation}
Using the standard values
$\zeta_H(0,q) = \tfrac{1}{2}-q$ and
$\zeta_H'(0,q) = \ln\Gamma(q) - \tfrac{1}{2}\ln(2\pi)$,
one computes:
\begin{align}
  \zeta(0;\mathcal{D}_a)
  &= \zeta_H(0,a)+\zeta_H(0,1-a)
   = \bigl(\tfrac{1}{2}-a\bigr)+\bigl(\tfrac{1}{2}-(1-a)\bigr) = 0,
  \label{eq:zeta0}\\[4pt]
  \zeta'(0;\mathcal{D}_a)
  &= \underbrace{\ln(R_\psi)\cdot\zeta(0;\mathcal{D}_a)}_{=0}
     + \bigl[\ln\Gamma(a)-\tfrac{1}{2}\ln(2\pi)\bigr]
     + \bigl[\ln\Gamma(1-a)-\tfrac{1}{2}\ln(2\pi)\bigr] \notag\\
  &= \ln\bigl[\Gamma(a)\,\Gamma(1-a)\bigr] - \ln(2\pi).
  \label{eq:zetaprime0}
\end{align}
Applying Euler's reflection formula $\Gamma(a)\,\Gamma(1-a) = \pi/\sin(\pi a)$:
\begin{equation}
  \zeta'(0;\mathcal{D}_a)
  = \ln\!\frac{\pi}{\sin(\pi a)} - \ln(2\pi)
  = -\ln\!\bigl(2\sin(\pi a)\bigr).
  \label{eq:zetaprime_final}
\end{equation}
Therefore
\begin{equation}
  \boxed{
  \det(\mathcal{D}_a)
  = \exp\bigl(-\zeta'(0;\mathcal{D}_a)\bigr)
  = 2\lvert\sin(\pi a)\rvert,
  \qquad a\notin\mathbb{Z}.
  }
  \label{eq:det_main}
\end{equation}
\textbf{This is $R_\psi$-independent.}
The $R_\psi^{-s}$ prefactor in~\eqref{eq:zeta_split} drops out because
$\zeta(0)=0$ (equation~\eqref{eq:zeta0}).
The result is a pure number determined solely by $a = \theta_H^*/(2\pi)$.

\subsection{Identification $Z^* = \det(\mathcal{D}_a)$}

The photon zero-mode functional integral contributes a factor
$(\det \mathcal{D}_a)^{-1/2}$ (bosonic field) or
$(\det \mathcal{D}_a)^{+1}$ (fermionic ghost) to the path integral.
The normalised zero-mode amplitude that enters the coupling is
\begin{equation}
  Z^* = \det(\partial_\psi + A_\psi) = 2\lvert\sin(\pi a)\rvert.
  \label{eq:Z_star_zetadet}
\end{equation}
With $\Lambda = 1/R_\psi$:
\begin{equation}
  \Phi_{\mathrm{UBT}}
  = \Lambda Z^*
  = \frac{2\lvert\sin(\pi a)\rvert}{R_\psi}.
  \label{eq:phi_L3_Rpsi}
\end{equation}
This still depends on $R_\psi$ through $\Lambda = 1/R_\psi$.
However, since $Z^*(R_\psi)$ is $R_\psi$-independent and $\Lambda \sim R_\psi^{-1}$,
the product $\Phi_{\mathrm{UBT}}$ scales as $1/R_\psi$, not as a pure number.

\subsection{Canonical vacuum $a=1$ ($\theta_H^* = 2\pi$): zero-mode problem}

For $\theta_H^* = 2\pi$ (i.e.\ $a=1$), the eigenvalue $\lambda_{-1} = 0$
is a \emph{zero mode} and $\det(\mathcal{D}_1) = 2\sin(\pi) = 0$.
The reduced determinant $\det'(\mathcal{D}_1)$ (excluding the zero mode)
is computed separately:

\noindent With $a=1$ and $n=-1$ excluded, the eigenvalues split as
$|n+1|$ for $n\geq 0$ (giving $\zeta_R(s)$) and for $n\leq -2$ (also
giving $\zeta_R(s)$), so:
\begin{equation}
  \zeta'(0;\mathcal{D}_1')
  = 2\zeta_R'(0) - \ln(R_\psi)\cdot 2\zeta_R(0)
  = 2\!\left(-\tfrac{1}{2}\ln(2\pi)\right) - \ln(R_\psi)\cdot(-1)
  = -\ln(2\pi) + \ln(R_\psi)
  = \ln\!\frac{R_\psi}{2\pi}.
\end{equation}
Therefore
\begin{equation}
  \det'(\mathcal{D}_1) = \frac{2\pi}{R_\psi},
  \label{eq:detprime_a1}
\end{equation}
and $\Phi_{\mathrm{UBT}} = \Lambda\,\det'(\mathcal{D}_1)
= (1/R_\psi)\cdot(2\pi/R_\psi) = 2\pi/R_\psi^2$,
which again depends on $R_\psi$.

\subsection{Numerical check}

For non-integer $a$ near $1$, equation~\eqref{eq:det_main} gives
$Z^* = 2\sin(\pi a)$.
The self-consistency equation~\eqref{eq:alpha_sc}
(with $1/(4\pi)$ prefactor and spectral term) combined
with $\Lambda Z^* = \Phi_{\mathrm{UBT}}$ requires
$\Phi_{\mathrm{UBT}} \approx 61$--$71$ (from the numerical analysis).
For $Z^*$ of order unity (e.g.\ $Z^* = 2\sin(\pi a)$ for generic $a$),
one needs $\Lambda \approx 61$--$71$, which means
$R_\psi \approx 1/65 \approx 0.015$.
This is far from the phenomenologically preferred $R_\psi \approx 0.917$,
confirming that the $Z^* = \det(\mathcal{D}_a)$ identification is
\textbf{not directly consistent} with the v1.1 equation at the level of
the algebraic part alone.

\begin{remark}
  The result~\eqref{eq:det_main} is nonetheless significant:
  it provides the \emph{only known $R_\psi$-independent quantity}
  derivable from $\mathcal{S}[\Theta]$ via zeta regularisation.
  If the definition $\Phi_{\mathrm{UBT}} = Z^* := 2|\sin(\pi a)|$ is adopted
  (without the $\Lambda$ factor), then the self-consistency equation
  fixes $R_\psi$ through $\Lambda = \Phi_{\mathrm{UBT}}/Z^* = 1$.
  This interpretation shifts the problem to explaining the unit value
  of $\Lambda$ in natural UBT units.
\end{remark}

\textbf{Conclusion for L3:}
The zeta-det approach yields the key result $\det(\mathcal{D}_a) = 2|\sin(\pi a)|$,
which is $R_\psi$-independent and topologically determined.
For the canonical vacuum $a=1$ a zero mode exists; $\det' = 2\pi/R_\psi$
still depends on $R_\psi$.
If $Z^* = 2|\sin(\pi a)|$ is the correct identification, the numerical
value $Z^* \leq 2$ is much smaller than the required $\Phi_{\mathrm{UBT}}\approx
61$--$71$ unless a large $\Lambda$ compensates.
Status: \textbf{[PARTIAL — key topological formula derived,
but numerical consistency with v1.1 equation not achieved without $R_\psi$ input]}.

% ─────────────────────────────────────────────────────────────────────────────
\section{Approach L4: Renormalization Finiteness Condition}
\label{sec:L4_renorm}
% ─────────────────────────────────────────────────────────────────────────────

\textbf{Status: [DEAD END — finiteness condition introduces renormalization
scale $\mu$ as a free parameter.]}

\subsection{Setup}

The photon self-energy (vacuum polarization) in 5D UBT at one loop is
\begin{equation}
  \Pi(q^2)
  = \frac{N_{\mathrm{eff}}}{12\pi^2}
    \Bigl[\Lambda^2 - q^2\ln\!\frac{\Lambda^2}{q^2} + \mathcal{O}(q^2)\Bigr].
  \label{eq:Pi_oneloop}
\end{equation}
The photon wavefunction renormalisation at scale $\mu$ is
\begin{equation}
  Z_3(\mu)
  = 1 - \frac{N_{\mathrm{eff}}}{12\pi^2}\ln\!\frac{\Lambda}{\mu}.
  \label{eq:Z3}
\end{equation}
The zero-mode norm from the renormalized propagator residue is
\begin{equation}
  Z^*(\mu)
  = \lim_{q^2\to 0}\frac{\Pi(q^2)}{q^2}
  = \frac{N_{\mathrm{eff}}}{12\pi^2}\ln\!\frac{\Lambda}{\mu}.
  \label{eq:Z_star_renorm}
\end{equation}

\subsection{Finiteness condition}

For $\Phi_{\mathrm{UBT}} = \Lambda Z^*(\mu)$ to be UV-finite
(independent of $\Lambda$ as $\Lambda\to\infty$), one requires
\begin{equation}
  Z^*(\mu) \;\sim\; \frac{c}{\Lambda}
  \quad\text{for some constant } c.
\end{equation}
From~\eqref{eq:Z_star_renorm} this demands
$\ln(\Lambda/\mu) \sim 1/\Lambda$, i.e.\
$\Lambda/\mu = e^{c'\!/\Lambda}\approx 1 + c'/\Lambda$,
which means $\mu \approx \Lambda$.
In other words, finiteness of $\Phi_{\mathrm{UBT}}$ requires the
renormalization scale $\mu$ to track the cutoff.

\subsection{Renormalization group consistency}

Setting $\mu = \mu_*$ such that $\Phi_{\mathrm{UBT}}(\mu_*) = \mathrm{const}$
gives the condition
\begin{equation}
  \frac{d}{d\mu}\Bigl[\Lambda\,Z^*(\mu)\Bigr] = 0
  \;\;\Longrightarrow\;\;
  \frac{d Z^*}{d\ln\mu} = 0
  \;\;\Longrightarrow\;\;
  -\frac{N_{\mathrm{eff}}}{12\pi^2} = 0.
  \label{eq:finiteness_impossible}
\end{equation}
This is impossible for $N_{\mathrm{eff}}\neq 0$.
The only escape is a cancellation between bosonic and fermionic
contributions such that the net $b_{\mathrm{eff}} = 0$,
which requires a supersymmetric spectrum not supported by the
$\mathbb{C}\otimes\mathbb{H}$ algebra with $N_{\mathrm{eff}} = 12$.

\textbf{Conclusion for L4:}
The renormalization finiteness condition does not uniquely fix $\mu$
or $\Phi_{\mathrm{UBT}}$; it merely relates $\mu$ to $\Lambda$ in a way
that reintroduces a free parameter.
Status: \textbf{[DEAD END]}.

% ─────────────────────────────────────────────────────────────────────────────
\section{Numerical Verification}
\label{sec:numerical}
% ─────────────────────────────────────────────────────────────────────────────

\subsection{$\alpha$ from self-consistency for each approach}

Using equation~\eqref{eq:alpha_sc} with
$\theta_H^{*2}/(4\pi(\mathcal{I}_\psi^*)^2) = 1/(4\pi)$,
$C_5 N_{\mathrm{eff}} = 3/(4\pi^2)$, and KK spectral contribution
$N_{\mathrm{eff}}/(2\pi)$ (equation~\eqref{eq:spectral_leading}):

\begin{center}
\begin{tabular}{lllll}
  \toprule
  Approach & $\Phi_{\mathrm{UBT}}$ expression
           & $R_\psi$ used
           & $\Phi_{\mathrm{UBT}}$
           & $\alpha^{-1}$ \\
  \midrule
  L1, canonical $Z^*$
    & $1/((2\pi)^3 R_\psi^2)$
    & $R_\psi=1$
    & $0.00405$
    & $\approx 445$ \\
  L1, algebraic trap
    & $1/((2\pi)^3 R_\psi^2)$
    & $R_\psi=\pi/\sqrt{12}$
    & $0.00496$
    & $\approx 370$ \\
  L1, fitted $R_\psi$
    & $1/((2\pi)^3 R_\psi^2)$
    & $R_\psi=0.917$
    & $0.00402$
    & $\approx 137$\,* \\
  L3, $Z^*=2|\sin(\pi a)|$
    & $2|\sin(\pi a)|/R_\psi$
    & $a\!\to\! 1, R_\psi\!=\!0.917$
    & $\to 0$
    & $\to\infty$ \\
  L3, $\det'(\mathcal{D}_1)$
    & $2\pi/R_\psi^2$
    & $R_\psi=0.917$
    & $7.46$
    & $\approx 25$ \\
  \bottomrule
\end{tabular}
\end{center}
\noindent
*\,This is the fitted value; $R_\psi = 0.917$ is \emph{not} derived here.

\medskip

The only approach yielding $\alpha^{-1} \approx 137$ is L1 with
$R_\psi \approx 0.917$ — but this is obtained by solving the self-consistency
equation backwards, not derived from first principles.
No approach yields $\alpha^{-1} = 137$ as an \emph{output}
without $R_\psi$ or $\alpha$ as an input.

% ─────────────────────────────────────────────────────────────────────────────
\section{Conclusion: Cumulative Status}
\label{sec:conclusion}
% ─────────────────────────────────────────────────────────────────────────────

\begin{center}
\begin{tabular}{lp{5.2cm}p{3.4cm}l}
  \toprule
  Approach & Key result & Blocker & Status \\
  \midrule
  L1 (zero-mode)
    & $\Phi_{\mathrm{UBT}} = 1/((2\pi)^3 R_\psi^2)$;\par
      trap: $R_\psi = \pi/\sqrt{N_{\mathrm{eff}}}$ (1\% accuracy)
    & $R_\psi$ not fixed algebraically
    & \textbf{[OPEN]} \\[6pt]
  L2 (Atiyah--Singer)
    & $\mathrm{ind}=0$ for canonical flat vacuum;\par
      twisted ($a\notin\mathbb{Z}$) case open
    & Integer holonomy forces trivial index
    & \textbf{[INCONCLUSIVE]} \\[6pt]
  L3 (zeta-det)
    & $\det(\partial_\psi+A_\psi) = 2|\sin(\pi a)|$;\par
      \emph{$R_\psi$-independent!}\par
      Canonical $a=1$: zero mode, $\det'=2\pi/R_\psi$
    & Numerical value $Z^*\leq 2$ too small;\par
      canonical vacuum is degenerate
    & \textbf{[PARTIAL]} \\[6pt]
  L4 (renorm.)
    & Finiteness condition~\eqref{eq:finiteness_impossible}
      is unsatisfiable
    & Introduces free scale $\mu$
    & \textbf{[DEAD END]} \\
  \bottomrule
\end{tabular}
\end{center}

\subsection*{Main finding of this document}

\begin{theorem}[Zeta-det formula]
  For the operator $\mathcal{D}_a = \partial_\psi + a/R_\psi$ on
  $S^1_\psi$ with $a = \theta_H^*/(2\pi)\notin\mathbb{Z}$,
  the $\zeta$-regularised functional determinant satisfies
  \begin{equation}
    \det(\mathcal{D}_a) = 2\lvert\sin(\pi a)\rvert,
    \label{eq:thm_zeta}
  \end{equation}
  independently of $R_\psi$.
  This is a pure number determined entirely by the Hosotani phase
  $\theta_H^*$.
\end{theorem}

\begin{proof}
  See computation in Section~\ref{subsec:L3_noninteger}:
  $\zeta(0;\mathcal{D}_a)=0$ eliminates all $R_\psi$ dependence,
  and $\zeta'(0;\mathcal{D}_a) = -\ln(2|\sin(\pi a)|)$ follows from the
  Hurwitz values $\zeta_H(0,q)=\frac{1}{2}-q$ and
  $\zeta_H'(0,q)=\ln\Gamma(q)-\frac{1}{2}\ln(2\pi)$
  combined with Euler's reflection formula.
\end{proof}

\subsection*{Open problem}

\begin{gap}
  \label{gap:phi_ubt_open}
  $\Phi_{\mathrm{UBT}}$ is not yet derived as a pure number from
  $\mathcal{S}[\Theta]$ alone.
  The canonical vacuum ($\theta_H^* = 2\pi$) has an integer holonomy
  that forces a zero mode, preventing a direct application of the
  $\det = 2|\sin(\pi a)|$ formula.
  Three sub-problems remain:
  \begin{enumerate}
    \item \emph{Fix $R_\psi$ algebraically.}
          The candidate $R_\psi = \pi/\sqrt{N_{\mathrm{eff}}}$
          (equation~\eqref{eq:R_psi_algebraic})
          is a \textbf{[NUMERICAL OBSERVATION]}; its derivation from
          $\mathcal{S}[\Theta]$ would close the gap.
    \item \emph{Regulate the zero mode in the canonical vacuum.}
          The zero mode of $\mathcal{D}_1$ must be handled separately
          (collective coordinate, modular space integral, or
          $\partial_\psi + a$ with $a\to 1^-$).
          A natural regulated value could recover $\Phi_{\mathrm{UBT}}$
          as a pure number.
    \item \emph{Link L2 (Atiyah--Singer) to L3.}
          If the index of the $\mathbb{C}\otimes\mathbb{H}$ bundle over
          a compact curved base is non-zero, it may set the overall
          normalisation of $\Phi_{\mathrm{UBT}}$ through the index theorem.
  \end{enumerate}
\end{gap}

\ifdefined\INCLUDEMODE\else
\end{document}
\fi
